\section{Discussion}

The principal empirical finding is that fusion strategy changes the value obtained from additional SAR inputs. The supplied test summaries show almost no benefit from early fusion, while separate encoders yield modestly lower errors. This pattern is consistent with a benefit from learning modality-specific features before fusion. It remains an interpretation: no feature-attribution experiment, capacity-matched model, or physical sensitivity analysis was conducted to establish the mechanism. The paper therefore treats E4 as the best observed configuration in this comparison, rather than a new general solution for optical--SAR learning.

The experimental controls strengthen the comparison by holding data assignment and training settings constant. However, multiple seeds on one fixed split assess training stochasticity rather than variability across geographic partitions. They cannot establish robustness to a different orientation of the boundaries, a new mangrove region, or a new acquisition period. Additional spatial folds and external references would be required to assess those questions. The 300 m buffers reduce direct boundary proximity but do not prove that spatial dependence has been eliminated.

Patch filtering is another part of the interpretation. Selecting patches with sufficient references and mangrove coverage creates a dataset suited to supervised training, but it also changes which environmental conditions are represented. The 118 excluded patches are not automatically an additional test set. Some can support dense inference if their predictors are valid; some fall in excluded buffers; others contain too few references to support evaluation. Any additional held-out assessment should be justified by region membership and available independent references, while preserving the original training boundary.

For practical mapping, the E4 mean test RMSE of 3.123 m remains above a 3 m target. The more appropriate claim is a reduction relative to the optical baseline, not attainment of that target. A 10 m grid is useful for alignment with other layers, but the accuracy evidence concerns agreement with GEDI RH98 at a small number of sampled cells. Combining the map with biomass equations would add uncertainties from allometry, forest composition, and reference support. This study supplies a structural prediction workflow and a bounded regional comparison; it does not validate a carbon inventory.
